\section{Introduction}\label{sec:introduction}

The congruence lattice of an algebra records all of its quotients at once.
The representation theorem of Gr\"atzer and Schmidt says that every
algebraic lattice occurs in this way \cite[Theorem~I, p.~34]{GratzerSchmidt1963}.
The finite version asks for a finite representing algebra whenever the
given lattice is finite. The difference is substantial: a finite
congruence lattice may be represented by an infinite algebra without
providing any finite representation.

Here a finite algebra is a finite nonempty set equipped with finitely many
total operations, each of finite arity. The signature is allowed to
depend on the lattice. Nullary operations and the empty signature are
allowed. We write $\Con(A)$ for the \emph{entire} lattice of congruence
relations on $A$, ordered by inclusion.
Allowing infinitely many finitary operations would give no additional
finite representations. Indeed, a finite carrier has only finitely many
equivalence relations. For each relation that is not a congruence, retain
one operation witnessing its failure of compatibility. The resulting
finite signature has exactly the same congruence lattice.

\begin{theorem}\label{main:congruences}
There exists a finite nonempty lattice $L$ such that
\[
 L\not\cong\Con(A)
\]
for every finite nonempty algebra $A$ of every finite signature.
\end{theorem}

Thus the finite lattice representation problem has a negative answer.
The theorem concerns unrestricted finite universal algebras; it is not
a restriction to a variety, a bound on the carrier size, or an obstruction
to a particular representation method.

P\'alfy and Pudl\'ak proved that the following two universal assertions
are equivalent \cite[Theorem~2]{PalfyPudlak1980}:
every finite lattice is the congruence lattice of a finite algebra, and
every finite lattice is an interval in the subgroup lattice of a finite
group. We establish the failure of the second assertion.

\begin{theorem}\label{main:group}
There exists a finite lattice that is not isomorphic to $[D,G]$ for any
subgroup $D$ of any finite group $G$.
\end{theorem}

Here $[D,G]=\{X:D\le X\le G\}$ carries subgroup inclusion. The equivalence
is global: Theorem~\ref{main:group} implies
Theorem~\ref{main:congruences}, without asserting that a prescribed
individual lattice has the two representation properties simultaneously.
The one-element lattice causes no exception; it is represented by a
one-element algebra and by a trivial subgroup interval.

The constructed subgroup-interval obstruction is self-dual.  It also
cannot be the entire lattice of intermediate fields of a finite separable
extension $E/F$.  Indeed, a
finite normal closure $N/F$ gives an order anti-isomorphism with
\[
 [\operatorname{Gal}(N/E),\operatorname{Gal}(N/F)],
\]
by the Galois correspondence; self-duality would turn such a field
representation into a forbidden subgroup-interval representation.
This is the group--field correspondence discussed in
\cite[Section~3]{PalfyTwisted}.

\subsection{History and the subgroup-interval approach}

The finite problem has several distinct formulations: existence of a
representation, finding a representing algebra, and deciding from a finite
order table whether a representation exists.  McNulty's account separates
these questions and places them in their universal-algebraic setting
\cite[pp.~29--30]{McNulty2015}.  We address the existence question by
ruling out a finite subgroup interval.  The global transfer just stated
is essential: the decorated lattice constructed below is a subgroup-interval
obstruction, and we do not identify it pointwise with the lattice asserted
by Theorem~\ref{main:congruences}.

Two positive representation results clarify the finite requirement.
Pudl\'ak and T\r{u}ma proved that every finite lattice embeds, preserving
meets and joins, into the partition lattice of a finite set
\cite{PudlakTuma1980}.  Such an embedding need not realize the given
lattice as the entire congruence lattice of operations on that set.
On the group side, Repnitski{\u{\i}} and T\r{u}ma represented every finite
lattice as an interval in the subgroup lattice of a countable locally
finite group \cite{RepnitskiiTuma2008}.  Local finiteness still permits
the ambient group to be infinite.

The subgroup-interval approach exposes normal subgroups and primitive
sections to lattice constraints.  Baddeley and Lucchini developed the
reduction of height-two interval questions to finite simple-group problems
\cite{BaddeleyLucchini1997}; Baddeley extended the class of lattice
configurations under consideration \cite{Baddeley1998}.
Further interval reductions were developed by B\"orner and Aschbacher
\cite{Borner1999,Aschbacher2008}.  These results have different hypotheses;
P\'alfy's comparison explains their relationship and the role of almost
simple groups and twisted wreath products \cite[Section~4]{PalfyTwisted}.
In particular, the invariant-subdirect and homomorphism-extension viewpoint
in \cite{PalfyTwisted} is the starting structure for our final reduction.
Aschbacher's minimum-order and signalizer-lattice viewpoint is a related
predecessor of the complementary-socle and extension-kernel arguments.

Lattice shapes can enforce group-theoretic conditions; DeMeo develops
this general perspective and the combination of such requirements
\cite{DeMeoInterval2014}.  Here the particular shapes must do two jobs
simultaneously.  They force uniform restrictions on every possible
simple-group-labelled chain, and they retain the Boolean joins on which
the final extensions must agree.  Uniformity is the obstacle: every
lattice in the construction is finite, but a hypothetical representing
group has no prescribed bound on its order, Lie rank, or field size.

\subsection{The proof in two stages}

Our construction has two stages. The first is a uniform rigidity result
for finite configurations of subgroups. We attach a nonabelian simple
type $T_X$ to selected vertices $X$.  If $D$ is the common bottom,
the unique minimal normal subgroup of $X/\operatorname{core}_X(D)$
is a direct power of $T_X$; here $\operatorname{core}_X(D)$ is the
largest normal subgroup of $X$ contained in $D$.
Special intervals, called fences, test normal subgroups.  Equal labels
give subdirect powers, in which selected simple coordinates are linked
diagonally.  At the comparisons used below, different labels give
coordinate products that can be projected to subgroup intervals of
almost simple groups.
Long chains are supplied with two-edge shortcuts and further interval
tests that make these projections available.  The crucial conclusion
is an \emph{absolute} upper bound
on the length of such a chain when its endpoint labels differ.
The bound must be independent of the finite group in which the
configuration is represented.

The classification of finite simple groups enters this stage through
maximal-subgroup structure and representation theory. A subgroup maximal
in a maximal subgroup of an almost simple group has at most five nonabelian chief
factors \cite[Proposition~8.1]{BurnessLiebeckShalev2017}.
We combine that bound with the fact that coordinate stabilizers in the
common bottom already induce the full local permutation actions on the
simple factors carrying the lower labels.
An elementary bounded-depth action argument then
controls the union of nonabelian simple types in the relevant projected
coordinate stabilizers, rather than just their number at one vertex. Decreasing
bottom cores account for the nonabelian layers lost between socles. These are the
uniformity statements needed before applying finite Ramsey theory.

The remaining classical-group comparisons are controlled by field
parameters, dimension growth, and irreducible restrictions.
The structural starting points are Aschbacher's maximal-subgroup
classification and its classical-group formulations
\cite{Aschbacher1984,KleidmanLiebeck1990}, together with Larsen--Pink's
fixed-type envelope theory \cite{LarsenPink2011}.
Landazuri--Seitz degree bounds separate cross-characteristic growth
\cite{LandazuriSeitz1974,Hasa2014}; Steinberg's representation theory and
L\"ubeck's restricted small-degree classification handle the defining
characteristic \cite{Steinberg1963,SteinbergNotes,Luebeck2001}.
The paper derives the required bounds on whole chains from these inputs;
none of these results alone bounds the lengths of the tested chains.
For defining-characteristic representations, the relevant action on
absolute constituent types is the difference between a source field
shift and an ambient semilinear shift. Their synchronization bounds
composition length independently of the field degree.
After another Ramsey argument, natural modules at later endpoints
restrict irreducibly to earlier ones.  The irreducible-triple input
comes from Seitz's classification and the corrections of Cavallin and
Testerman, in the formulation of Burness and Testerman
\cite{Seitz1987,CavallinTesterman2019,BurnessTesterman2018}.
We use it twice: first with a bound on the composition length of the
ambient natural module on the smallest subgroup
(Lemma~\ref{irr:rank-length}), and then with an intervening simple
subgroup of strictly larger rank whose natural module is irreducible
on that smallest subgroup (Lemma~\ref{irr:intervening}).
These hypotheses control the twisted-diagonal spin families.
The resulting rank bounds are incompatible with a sufficiently long
chain and its two-edge shortcuts.

In the second stage, we build one finite lattice from a Boolean lattice
of rank four and a small additional configuration, called the detector.
The chain tests are inserted in both required orders while preserving
the original Boolean meets and joins.  Their length is chosen from the
absolute bound, before any representing group is considered.

Suppose the resulting lattice is represented by $[D,G]$, with $|G|$
minimal among representations in either orientation.  The global fence
gives a socle $S\cong T^I$ and $G=DS$; minimality gives $D\cap S=1$.
Equality of the forward labels forces every proper nonzero original
Boolean vertex $Y$ to meet $S$ in a subdirect subgroup, meaning one
which projects onto every simple coordinate.  Fixing a coordinate
gives its stabilizer $A\le D$ and an action
\[
                   \alpha:A\longrightarrow\Aut(T).
\]
A $D$-invariant subdirect subgroup corresponds to a subgroup $U$ with
$A\le U\le D$ and a homomorphism $\beta:U\to\Aut(T)$ extending
$\alpha$.  This correspondence reverses inclusion.  It realizes the
reversed upper filter above $Y$ as the actual group interval $[A,U_Y]$.

The reverse chain tests make the socles of
$U_Y/\operatorname{core}_{U_Y}(A)$ powers of one simple group $R$,
which need not equal $T$.
The detector ensures $\alpha(A)\ge\Inn(T)$.  Read the original
Boolean copy in reverse order.  Proper Boolean unions give extensions
on generated domains that agree on their overlaps.  For a full union,
the corresponding forward vertices meet at $D$, so their intersections
with $S$ have trivial intersection.  A common extension would produce
a nontrivial subdirect subgroup in this intersection and is therefore
forbidden.

The final argument forces that forbidden extension.  If the core of
$A$ in an atom domain had trivial image, compatible lifts of socle
factors on the three rank-two faces through that atom would already
produce a common extension.  Hence every atom core has image containing
$\Inn(T)$.  Intersections give the same
property for larger proper faces.  A least normal subgroup of $A$
with this inner image is then normal in all the relevant domains;
conjugation on its quotient isomorphic to $T$ supplies their common
extension.  This is the contradiction.

\subsection{Structure and scope of the proof}

Section~\ref{sec:framework} develops the lattice tests and factor
transport. The central group-theoretic result is
Theorem~\ref{cls:chain-rigidity}, stated in
Section~\ref{sec:classification} and proved in
Sections~\ref{sec:classification}--\ref{sec:rigidity}.
These sections control, successively, factor actions and simple types,
field and dimension growth, algebraic lifts along paths,
natural-module composition length, and irreducible restrictions.
Section~\ref{sec:lattices} constructs a finite lattice realizing all
the tests. Sections~\ref{ext:section} and~\ref{glue:section} reduce a
hypothetical representation to an extension configuration and rule
it out by Boolean gluing, proving Theorem~\ref{main:group}.

The finite lattice is specified by finitely many private-branch
insertions with a sufficiently large finite chain parameter.
An explicit numerical value of that parameter is unnecessary for the
existence conclusion. Its existence follows from bounds that are
uniform over all possible finite-group representations.

All groups are finite unless called algebraic groups. Algebraic
groups occur only in the representation-theoretic part of the proof
and are taken over algebraically closed fields of positive
characteristic, with central covers and the stated special isogenies
when needed. We use the classification of finite simple groups and
identify each additional classification or representation input at
its point of application. The new uniformity, lifting, lattice
bookkeeping and gluing deductions are proved explicitly.
